% !TEX spellcheck = en_US
% !TEX spellcheck = LaTeX
\documentclass[letterpaper,10pt,english]{article}
\usepackage{%
amsfonts,%
amsmath,%
amssymb,%
amsthm,%
babel,%
bbm,%
%biblatex,%
caption,%
centernot,%
color,%
enumerate,%
%enumitem,%
epsfig,%
epstopdf,%
etex,%
fancybox,%
framed,%
fullpage,%
%geometry,%
graphicx,%
hyperref,%
latexsym,%
mathptmx,%
mathtools,%
multicol,%
paralist,%
pgf,%
pgfplots,%
pgfplotstable,%
pgfpages,%
proof,%
psfrag,%
%subfigure,%
tcolorbox,%
tfrupee,%
tikz,%
times,%
%ulem,%
url,%
xcolor,%
mathpazo,%
}

\definecolor{shadecolor}{gray}{.95}%{rgb}{1,0,0}
\usepackage[margin=1in,top=0.75in]{geometry}
\usepackage[mathscr]{eucal}
\usepgflibrary{shapes}
\usepgfplotslibrary{fillbetween}
\usetikzlibrary{%
arrows,%
backgrounds,%
chains,%
decorations.pathmorphing,% /pgf/decoration/random steps | erste Graphik
decorations.text,% 
matrix,%
positioning,% wg. " of "
fit,%
patterns,%
petri,%
plotmarks,%
scopes,%
shadows,%
shapes.misc,% wg. rounded rectangle
shapes.arrows,%
shapes.callouts,%
shapes%
}

%\pgfplotsset{compat=newest} %<------ Here
\pgfplotsset{compat=1.11} %<------ Or use this one

\theoremstyle{plain}
\newtheorem{thm}{Theorem}[section]
\newtheorem{lem}[thm]{Lemma}
\newtheorem{prop}[thm]{Proposition}
\newtheorem{cor}[thm]{Corollary}
\newtheorem{clm}[thm]{Claim}

\theoremstyle{definition}
\newtheorem{defn}[thm]{Definition}
\newtheorem{conj}[thm]{Conjecture}
\newtheorem{exmp}[thm]{Example}
\newtheorem{axiom}[thm]{Axiom}
\newtheorem{assum}[thm]{Assumption}
\newtheorem{exerc}[thm]{Exercise}
\newtheorem*{defin}{Definition}

\theoremstyle{remark}
\newtheorem{rem}{Remark}
\newtheorem{note}{Note}

\newcommand{\iid}{\emph{i.i.d.}\ }
\newcommand{\Cov}{\operatorname{Cov}}
%\newcommand{\det}{\operatorname{det}}
%\newcommand{\Real}{\mathbb{R}}
\newcommand{\tr}{\operatorname{tr}}
\newcommand{\Var}{\operatorname{Var}}
\newcommand{\cov}{\operatorname{cov}}

%\renewcommand{\proof}[1]{\begin{proof}#1\end{proof}}
\newcommand{\EQ}[1]{\begin{equation*}#1\end{equation*}}
\newcommand{\EQN}[1]{\begin{equation}#1\end{equation}}
\newcommand{\eq}[1]{\begin{align*}#1\end{align*}}
\newcommand{\meq}[2]{\begin{xalignat*}{#1}#2\end{xalignat*}}
\newcommand{\norm}[1]{\left\lVert#1\right\rVert}
\newcommand{\inner}[1]{\left\langle #1\right\rangle}
\newcommand{\abs}[1]{\left\lvert#1\right\rvert}
\newcommand{\expect}[1]{\mathbb{E}\left[{#1}\right]}
\newcommand{\prob}[1]{\mathbb{P}\left[{#1}\right]}
\newcommand{\given}{\; \big\vert \;} 
\newcommand{\set}[1]{\left\{#1\right\}} 
\newcommand{\indicator}[1]{1_{\set{#1}}} 
\newcommand{\Ind}[1]{\mathbbm{1}_{#1}} 
\newcommand{\SetIn}[1]{\mathbbm{1}_{\set{#1}}} 
\newcommand{\red}[1]{\textcolor{red}{#1}} 
\newcommand{\floor}[1]{\left\lfloor#1\right\rfloor}
\newcommand{\ceil}[1]{\left\lceil#1\right\rceil}


\newcommand{\C}{\mathbb{C}}
\newcommand{\D}{\mathbb{D}}
\newcommand{\E}{\mathbb{E}}
\newcommand{\F}{\mathbb{F}}
\newcommand{\N}{\mathbb{N}}
\renewcommand{\P}{\mathbb{P}}
\newcommand{\Q}{\mathbb{Q}}
\newcommand{\R}{\mathbb{R}}
\newcommand{\Z}{\mathbb{Z}}

\newcommand{\bA}{\mathbf{A}}
\newcommand{\bI}{\mathbf{I}}
\newcommand{\bU}{\mathbf{1}}
\newcommand{\bx}{\mathbf{x}}
\newcommand{\bX}{\mathbf{X}}
\newcommand{\bZ}{\mathbf{Z}}


\newcommand{\cA}{\mathcal{A}}
\newcommand{\cB}{\mathcal{B}}
\newcommand{\cC}{\mathcal{C}}
\newcommand{\cF}{\mathcal{F}}
\newcommand{\cG}{\mathcal{G}}
\newcommand{\cM}{\mathcal{M}}
\newcommand{\cN}{\mathcal{N}}
\newcommand{\cP}{\mathcal{P}}
\newcommand{\cT}{\mathcal{T}}
\newcommand{\cX}{\mathcal{X}}
\newcommand{\cY}{\mathcal{Y}}

\newcommand{\sA}{\mathscr{A}}
\newcommand{\sB}{\mathscr{B}}
\newcommand{\sC}{\mathscr{C}}
\newcommand{\sE}{\mathscr{E}}
\newcommand{\sF}{\mathscr{F}}
\newcommand{\sG}{\mathscr{G}}
\newcommand{\sH}{\mathscr{H}}
\newcommand{\sL}{\mathscr{L}}
\newcommand{\sS}{\mathscr{S}}
\newcommand{\sT}{\mathscr{T}}
\newcommand{\sX}{\mathscr{X}}
\newcommand{\sY}{\mathscr{Y}}
\newcommand{\sZ}{\mathscr{Z}}

\newcommand{\tS}{\tilde{S}}
% Debug
\newcommand{\todo}[1]{\begin{color}{blue}{{\bf~[TODO:~#1]}}\end{color}}

% a few handy macros

\renewcommand{\le}{\leqslant}
\renewcommand{\ge}{\geqslant}
\newcommand\matlab{{\sc matlab}}
\newcommand{\goto}{\rightarrow}
\newcommand{\bigo}{{\mathcal O}}
%\newcommand{\half}{\frac{1}{2}}
%\newcommand\implies{\quad\Longrightarrow\quad}
\newcommand\reals{{{\rm l} \kern -.15em {\rm R} }}
\newcommand\complex{{\raisebox{.043ex}{\rule{0.07em}{1.56ex}} \hskip -.35em {\rm C}}}


% macros for matrices/vectors:

% matrix environment for vectors or matrices where elements are centered
\newenvironment{mat}{\left[\begin{array}{ccccccccccccccc}}{\end{array}\right]}
\newcommand\bcm{\begin{mat}}
\newcommand\ecm{\end{mat}}

% matrix environment for vectors or matrices where elements are right justifvied
\newenvironment{rmat}{\left[\begin{array}{rrrrrrrrrrrrr}}{\end{array}\right]}
\newcommand\brm{\begin{rmat}}
\newcommand\erm{\end{rmat}}

% for left brace and a set of choices
%\newenvironment{choices}{\left\{ \begin{array}{ll}}{\end{array}\right.}
\newcommand\when{&\text{if~}}
\newcommand\otherwise{&\text{otherwise}}
% sample usage:
%\delta_{ij} = \begin{choices} 1 \when i=j, \\ 0 \otherwise \end{choices}


% for labeling and referencing equations:
\newcommand{\eql}{\begin{equation}\label}
\newcommand{\eqn}[1]{(\ref{#1})}
% can then do
%\eql{eqnlabel}
%...
%\end{equation}
% and refer to it as equation \eqn{eqnlabel}.
% some useful macros for finite difference methods:
\newcommand\unp{U^{n+1}}
\newcommand\unm{U^{n-1}}
% for chemical reactions:
\newcommand{\react}[1]{\stackrel{K_{#1}}{\rightarrow}}
\newcommand{\reactb}[2]{\stackrel{K_{#1}}{~\stackrel{\rightleftharpoons}
 {\scriptstyle K_{#2}}}~}
\makeatletter
\def\th@plain{%
\thm@notefont{}% same as heading font
\itshape % body font
}
\def\th@definition{%
\thm@notefont{}% same as heading font
\normalfont % body font
}
\makeatother
\date{}


\title{Lecture-14: Almost sure convergence}% of random variables}
\author{}

\begin{document}
\maketitle

\section{Point-wise convergence}
Consider a random sequence $X:\Omega\to\R^\N$ defined on a probability space $(\Omega, \sF, P)$, 
then each $X_n \triangleq \pi_n\circ X: \Omega \to \R$ is a random variable. 
%Recall that a random variable $X:\Omega \to \R$, is an $\sF$-measurable function on the sample space $\Omega$ such that $X^{-1}(-\infty, x] \in \sF$ for all $x \in \R$. 
%A sequence of random variables $X:\Omega\to\R^\N$ is hence a sequence of $\sF$-measurable functions. 
There are many possible definitions for convergence of a sequence of random variables. 
One idea is to consider $X(\omega)\in \R^\N$ as a real valued sequence for each outcome $\omega$, 
and consider the $\lim_nX_n(\omega)$ for each outcome $\omega$. 
\begin{defn}%[point wise convergence] 
A random sequence $X:\Omega\to\R^\N$ defined on a probability space $(\Omega, \sF, P)$ \textbf{converges point-wise} to a random variable $X_\infty:\Omega \to \R$, if for all outcomes $\omega \in \Omega$, we have 
\EQ{
\lim_nX_n(\omega) = X_\infty(\omega).
}
\end{defn}
\begin{rem} 
This is a very strong convergence. 
Intuitively, what happens on an event of probability zero is not important. 
We will strive for a weaker notion of convergence, 
where the sequence of random variable converge point-wise on a set of outcomes with probability one.
\end{rem}


\section{Almost sure statements}
\begin{defn} 
A statement holds \textbf{almost surely} (a.s.) if there exists an event called the \textbf{exception set} $N \in \sF$ with $P(N) = 0$ such that the statement holds for all $\omega \notin N$. 
%The set $N$ is called the \textbf{exception set}. 
\end{defn}
\begin{shaded}
\begin{exmp}[Almost sure equality] 
Two random variables $X, Y$ defined on the probability space $(\Omega, \sF, P)$ are said to be equal a.s. if the following exception set 
\EQ{
N \triangleq \set{\omega \in \Omega: X(\omega) \neq Y(\omega)} \in \sF,
}
has probability measure $P(N) = 0$. 
Then $Y$ is called a \textbf{version} of $X$, and we can define an equivalence class of a.s. equal random variables. 
\end{exmp}
%\end{shaded}
%\begin{shaded}
\begin{exmp}[Almost sure monotonicity] 
Two random variables $X, Y$ defined on the probability space $(\Omega, \sF, P)$ are said to be $X \le Y$ a.s. if the exception set $N \triangleq \set{\omega \in \Omega: X(\omega) > Y(\omega)} \in \sF$ has probability measure $P(N) = 0$. 
\end{exmp}
\end{shaded}

\section{Almost sure convergence}

\begin{defn}[Almost sure convergence]
A random sequence $X:\Omega\to\R^\N$ defined on the probability space $(\Omega, \sF, P)$ \textbf{converges almost surely}, 
if the following exception set 
\EQ{
N \triangleq \set{\omega \in \Omega:  \lim\inf_nX_n(\omega) < \lim\sup_nX_n(\omega) \text{ or } \lim\sup_nX_n(\omega) = \infty} \in \sF,
}
has zero probability. 
%then $\lim_n X_n$ exists a.s. means there exists an exception event $N \in \sF$, 
%such that $P(N) = 0$ and if $\omega\notin N$, then $\lim_{n}X_n(\omega)$ exists. 
%That is, 
%\EQ{
%N^c = \set{\omega \in \Omega: \lim\sup_nX_n(\omega) = \lim\inf_nX_n(\omega)}. 
%}
Let $X_\infty$ be the point-wise limit of the sequence of random variables $X:\Omega\to\R^\N$ on the set $N^c$, 
then we say that the sequence $X$ \textbf{converges almost surely} to $X_\infty$, and denote it as 
\EQ{
\lim_nX_n = X_\infty \text{ a.s.}
}
\end{defn}
\begin{shaded}
\begin{exmp}[Convergence almost surely but not everywhere] 
Consider the probability space $([0,1], \sB([0,1]), \lambda)$ such that $\lambda([a,b]) = b-a$ for all $0 \le a \le b \le 1$. 
For each $n \in \N$, we define the scaled indicator random variable $X_n: \Omega \to \set{0,1}$ such that
\EQ{
X_n(\omega) \triangleq n\Ind{\left[0, \frac{1}{n}\right]}(\omega).
}
Let $N = \set{0}$, then for any $\omega \notin N$, there exists $m = \lceil\frac{1}{\omega}\rceil \in \N$, such that for all $n > m$, we have $X_n(\omega) = 0$. 
That is, $\lim_nX_n = 0$ a.s. since $\lambda(N) = 0$. 
However, $X_n(0) = n$ for all $n \in \N$. 
\end{exmp}
\end{shaded}

\section{Convergence in probability}
\begin{defn}[convergence in probability] 
A random sequence $X:\Omega\to\R^\N$ defined on the probability space $(\Omega, \sF, P)$ \textbf{converges in probability} to a random variable $X_\infty:\Omega\to\R$, 
if $\lim_nP(A_n(\epsilon)) = 0$ for any $\epsilon > 0$, 
where 
\EQ{
A_n(\epsilon) \triangleq \set{\omega \in \Omega: \abs{X_n(\omega)-X_\infty(\omega)} > \epsilon} \in \sF.
}
\end{defn}
\begin{rem} 
%For a sequence $X:\Omega\to\R^\N$, 
$\lim_nX_n = X_\infty$ a.s. means that for almost all outcomes $\omega$, 
the difference $X_n(\omega)-X_\infty(\omega)$ gets small and stays small. 
\end{rem}
\begin{rem}
%On the other hand, %convergence in probability
$\lim_nX_n = X_\infty$ i.p. is a weaker convergence than a.s. convergence, 
and merely requires that the probability of the difference $X_n(\omega)-X_\infty(\omega)$ being non-trivial becomes small. 
\end{rem}
\begin{shaded}
\begin{exmp}[Convergence in probability but not almost surely] 
Consider the probability space $([0,1], \sB([0,1]), \lambda)$ such that $\lambda([a,b]) = b-a$ for all $0 \le a \le b \le 1$. 
For each $k \in \N$, we consider the sequence $S_k = \sum_{i=1}^ki$, and define integer intervals $I_k \triangleq \set{S_{k-1}+1, \dots, S_k}$. 
Clearly, the intervals $(I_k:k \in \N)$ partition the natural numbers where $\abs{I_k} = k$. 
It follows that each $n \in \N$ lies in some $I_{k_n}$, such that $n = S_{k_n-1}+i_n$ for $i_n \in [k_n]$. 
Therefore, for each $n \in \N$, we define indicator random variable $X_n: \Omega \to \set{0,1}$ such that 
\EQ{
X_n(\omega) = \Ind{\left[\frac{i_n-1}{k_n}, \frac{i_n}{k_n}\right]}(\omega).
}
For any $\omega \in [0,1]$, we have $X_n(\omega) = 1$ for infinitely many values since there exist infinitely many $(i,k)$ pairs such that $\frac{(i-1)}{k} \le \omega \le \frac{i}{k}$, 
and hence $\lim\sup_nX_n(\omega) = 1$ and hence $\lim_nX_n(\omega) \neq 0$. 
However, $\lim_nX_n(\omega) = 0$ in probability, since 
\EQ{
\lim_n\lambda\set{X_n(\omega) \neq 0} = \lim_n\frac{1}{k_n} = 0.
}
\end{exmp}
\end{shaded}

\section{Infinitely often and all but finitely many}
\begin{lem}[infinitely often and all but finitely many]
Let $A \in \sF^\N$ be a sequence of events. 
\begin{enumerate}[(a)]
\item For some subsequence $(k_n: n \in \N)$ depending on $\omega$, we have 
\EQ{
\lim\sup_nA_n 
= \set{\omega \in \Omega: \omega \in A_{k_n} \text{ for all } n \in \N} 
= \set{\omega \in \Omega: \sum_{n \in \N}\Ind{A_n}(\omega) = \infty} 
= \set{A_n \text{ infinitely often }}.
}
\item For a finite $n_0(\omega) \in \N$ depending on $\omega$, we have 
\EQ{
\lim\inf_nA_n = \set{\omega \in \Omega: \omega \in A_n \text{ for all  }n \ge n_0(\omega)}
= \set{\omega \in \Omega: \sum_{n \in \N}\Ind{A_n^c}(\omega) < \infty}
= \set{A_n \text{ for all but finitely many }n}.
%= \set{\omega \in \Omega: \omega \in A_{n_k}, k \in \N} = \set{A_n \text{infinitely often }}.
}
\end{enumerate}
\end{lem}
\begin{proof} 
Let $A \in \sF^\N$ be a sequence of events. 
\begin{enumerate}[(a)]
\item Let $\omega \in \lim\sup_nA_n = \cap_{n \in \N}\cup_{k \ge n}A_k$, then $\omega \in \cup_{k \ge n}A_k$ for all $n \in \N$. 
Therefore, for each $n \in \N$, there exists $k_n \in \N$ such that $\omega \in A_{k_n}$, and hence  
\EQ{
\sum_{j  \in \N}\Ind{A_j}(\omega) 
\ge \sum_{n  \in \N}\Ind{A_{k_n}}(\omega) = \infty. 
}
Conversely, if $\sum_{j  \in \N}\Ind{A_j}(\omega) = \infty$, then for each $n \in \N$ there exists a $k_n \in \N$ such that $\omega \in A_{k_n}$ and hence $\omega \in \cup_{k \ge n}A_k$ for all $n \in \N$.

\item Let $\omega\in \lim\inf_nA_n = \cup_{n\in\N}\cap_{k\ge n}A_k$, then there exists $n_0(\omega)$ such that $\omega \in A_n$ for all $n \ge n_0(\omega)$. 
Conversely, if $\sum_{j  \in \N}\Ind{A_j^c}(\omega) < \infty$, then there exists $n_0(\omega)$ such that $\omega \in A_n$ for all $n \ge n_0(\omega)$.
\end{enumerate}
\end{proof}

\begin{thm}[Convergence a.s. implies in probability] 
If a sequence of random variables $X: \Omega \to \R^\N$ defined on a probability space $(\Omega, \sF, P)$ converges a.s. to a random variable $X_\infty:\Omega\to\R$, then it converges in probability to the same random variable.  
\end{thm}
\begin{proof}  
Let $\lim_nX_n = X_\infty$ a.s. and $\epsilon > 0$. 
We define events $A_n \triangleq \set{\omega \in \Omega: \abs{X_n(\omega)-X_\infty(\omega)} > \epsilon}$ for each $n \in \N$. 
We will show that %if $\lim_nX_n = X_\infty$ a.s., then 
$\lim_nP(A_n) = 0$. 
To this end, let $N$ be the exception set such that 
\EQ{
N \triangleq \set{\omega\in\Omega: \lim\inf_nX_n(\omega) < \lim\sup_nX_n(\omega) \text{ or } \lim\sup_nX_n(\omega) = \infty}.
}  
For $\omega \notin N$, there exists an $n_0(\omega)$ such that $\abs{X_n-X_\infty} \le  \epsilon$ for all $n \ge n_0$. 
That is, $\omega \in A_n^c$ for all $n \ge n_0(\omega)$ and hence $N^c \subseteq \lim\inf_nA_n^c$. 
It follows that $1 = P(\lim\inf_nA_n^c)$.   
%\EQ{
%1 = P(N^c) \le P(A_n^c \text{ for all but finitely many }n) = P(\lim\inf_nA_n^c) \le 1. 
%}
Since  $\lim\inf_nA_n^c = (\lim\sup_nA_n)^c$, we get 
%\EQ{
$0 = P(\lim\sup_nA_n) = \lim_nP(\cup_{k \ge n}A_k) \ge \lim_nP(A_n) \ge 0.$
%}
\end{proof}
\section{Borel-Cantelli Lemma}
\begin{prop}[Borel-Cantelli Lemma] 
Let $A \in \sF^\N$ be a sequence of events such that $\sum_{n \in \N}P(A_n) < \infty$, 
then $P\set{A_n \text{ i.o.}} = 0$. 
\end{prop}
\begin{proof} 
We can write the probability of infinitely often occurrence of $A_n$, by the continuity and sub-additivity of probability as  
\EQ{
P(\lim\sup_nA_n) = \lim_nP(\cup_{k\ge n}A_k) \le \lim_n\sum_{k \ge n}P(A_k) = 0.
}
The last equality follows from the fact that $\sum_{n \in \N}P(A_n) < \infty$. 
\end{proof}
\begin{prop}[Borel zero-one law]
Let $A \in \sF^\N$ be a sequence of independent events, then
\EQ{
P\set{A_n \text{ i.o.}} = \begin{cases}
0, & \text{ iff } \sum_nP(A_n) < \infty,\\
1, & \text{ iff } \sum_nP(A_n) = \infty.
\end{cases}
}
\end{prop}
\begin{proof} 
Let $A \in \sF^\N$ be a sequence of independent events. 
\begin{enumerate}[(a)]
\item From Borel-Cantelli Lemma, if $\sum_nP(A_n) < \infty$ then $P\set{A_n\text{ i.o.}} = 0$. 
\item Conversely, suppose $\sum_nP(A_n) = \infty$, 
then $\sum_{k \ge n}P(A_k) = \infty$ for all $n \in \N$.  
From the definition of $\lim\sup$ and $\lim\inf$, continuity of probability, 
and independence of sequence of events $A \in \sF^\N$, we get  
\EQ{
P\set{A_n\text{ i.o.}} = 1 - P(\lim\inf_nA_n^c) = 1 - \lim_n\lim_mP(\cap_{k=n}^mA_k^c) = 1 - \lim_n\lim_m\prod_{k=n}^m(1-P(A_k)).
}
Since $1-x \le e^{-x}$ for all $x \in \R$, from the above equation, the continuity of exponential function, and the hypothesis, we get 
\EQ{
1 \ge P\set{A_n\text{ i.o.}} 
\ge 1 -  \lim_n\lim_me^{-\sum_{k=n}^mP(A_k)} 
= 1 - \lim_n\exp(-\sum_{k\ge n}P(A_k)) 
= 1.
}
\end{enumerate}
\end{proof}
\begin{shaded}
\begin{exmp}[Convergence in probability can imply almost sure convergence] 
Consider a random Bernoulli sequence $X: \Omega \to \set{0,1}^\N$ defined on the probability space $(\Omega, \sF, P)$ such that $P\set{X_n = 1} = p_n$ for all $n \in \N$. 
Note that the sequence of random variables is not assumed to be independent, and definitely not identical. 
If $\lim_np_n = 0$, then we see that $\lim_nX_n = 0$ in probability. 

In addition, if $\sum_{n\in \N}p_n <\infty$, then $\lim_nX_n = 0$ a.s. 
To see this, we define event $A_n \triangleq \set{X_n = 1} \in \sF$ for each $n \in \N$. 
Then, applying the Borel-Cantelli Lemma to sequence of events $A \in \sF^\N$, 
we get  
\EQ{
1 = P((\lim\sup_nA_n)^c) = P(\lim\inf_nA_n^c).
}
That is, $\lim_nX_n = 0$ for $\omega \in \lim\inf_nA_n^c$, implying almost sure convergence. 
\end{exmp}
\end{shaded}
\red{
\begin{thm} 
A random sequence $X: \Omega\to\R^\N$ converges to a random variable $X_\infty:\Omega\to\R$ in probability, 
then there exists a subsequence $({n_k}: k \in \N) \subset \N$ such that $(X_{n_k}: k \in \N)$ converges almost surely to $X_\infty$.
\end{thm}
\begin{proof}
Letting $n_1 = 1$, we define the following subsequence and event recursively for each $j \in \N$, 
\meq{2}{
&n_j \triangleq \inf\set{N > n_{j-1}: P\set{\abs{X_r-X_\infty}>2^{-j}}<2^{-j}, \text{ for all } r \ge N}, &
&A_j \triangleq \set{\abs{X_{n_{j+1}}-X_\infty}>2^{-j}}.
}
From the construction, we have $\lim_kn_k = \infty$, and $P(A_j)<2^{-j}$ for each $j \in \N$. 
Therefore, $\sum_{k \in \N}P(A_k) < \infty$, 
and hence by the Borel-Cantelli Lemma, we have $P(\lim\sup_kA_{k}) = 0$.
Let $N = \lim\sup_kA_k$ be the exception set such that for any outcome $\omega \notin N$, for all but finitely many $j \in \N$
\EQ{
\abs{X_{n_j}(\omega)-X_\infty(\omega)} \le 2^{-j}.
}
That is, for all $\omega \notin N$, we have $\lim_nX_n(\omega) = X_\infty(\omega)$. 
\end{proof}
\begin{thm} 
A random sequence $X:\Omega\to\R^\N$ converges to a random variable $X_\infty$ in probability 
iff each subsequence $(X_{n_k}: k \in \N)$ contains a further subsequence $(X_{n_{k_j}}: j \in \N)$ converging almost surely to $X_\infty$.
\end{thm}
}
\appendix
\section{Limits of sequences}
\begin{defn} 
For any real valued sequence $a \in \R^\N$, we can define
\begin{xalignat*}{2}
&\lim\sup_n a_n \triangleq \inf_n\sup_{k \ge n}a_k,&
&\lim\inf_n a_n \triangleq \sup_n\inf_{k \ge n}a_k.
\end{xalignat*}
\end{defn}
\begin{rem} 
We define $e_n \triangleq \sup_{k \ge n}a_k$ and $f_n \triangleq \inf_{k \ge n}a_k$, and observe that $f_n \le a_k$ for all $k \ge n$. 
That is, $f_1, \dots, f_{n-1} \le a_n$ and $f_k \le a_k$ for all $k \ge n$. 
It follows that $\sup_nf_n \le \sup_{k \ge n}a_k  = e_n$ for all $n \in \N$, and hence $\lim\inf_n a_n  = \sup_nf_n \le \inf_n e_n =  \lim\sup_na_n$. 
\end{rem}
\begin{defn} 
A sequence $a \in \R^\N$ is said to converge if $\lim\sup_n a_n = \lim\inf_n a_n$ and the limit is defined as $a_n \triangleq \lim_n a_n = \lim\sup_n a_n = \lim\inf_n a_n. $
\end{defn}
\begin{thm} 
A sequence $a \in \R^\N$ converges to $a_\infty \in \R$ if for all $\epsilon > 0$ there exists an integer $N \in \N$ such that for all $n \ge N$, we have $\abs{a_n-a_\infty} < \epsilon$. 
\end{thm}
\begin{proof} 
Let $\epsilon > 0$ and find the integer $N_\epsilon \in \N$ such that $a_n \in (a_\infty-\epsilon, a_\infty+\epsilon)$ for all $n \ge N_\epsilon$. 
It follows that $a_\infty-\epsilon \le f_n \le e_n \le a_\infty+\epsilon$ for all $n \ge N_\epsilon$, 
and hence $a_\infty-\epsilon \le \lim\inf_n a_n \le \lim\sup_n a_n \le a_\infty+\epsilon.$ 
Since $\epsilon$ was arbitrary, it follows that $\lim_na_n = a_\infty$.
\end{proof}
\begin{prop}
For any sequence $a \in \R_+^\N$, the following statements are true.
\begin{compactenum}[(i)]
\item If $\sum_{n\in\N}a_n < \infty$ then $\lim_{n \to \infty}\sum_{k \ge n}a_k = 0$.
\item If $\sum_{n\in\N}a_n = \infty$ then $\sum_{k \ge n}a_k = \infty$ for all $k \in \N$.
\end{compactenum}
\end{prop}
\begin{proof}
We observe that $(\sum_{k < n}a_k: n \in\N)$ is a non-decreasing sequence, and hence $\lim_{n \to \infty}\sum_{k < n}a_k = \sup_n\sum_{k< n}a_k = \sum_{n \in \N}a_n$. 
\begin{compactenum}[(i)] 
\item It follows that $\sum_{k \ge n}a_k = \sum_{n \in \N}a_n - \sum_{k < n}a_k$ is a non-increasing sequence with limit $0$. 
\item We can write $\sum_{n \in \N}a_n = \sum_{k < n}a_k + \sum_{k \ge n}a_k$. 
Since the first term is finite for all $n \in \N$, it follows that the second term must be infinite for all $n \in \N$. 
\end{compactenum}
\end{proof}
\end{document}